\chapter{Topologia gładkich zespolonych odmian projektowych: twierdzenia Lefschetza, przecięcia i rozdmuchania}
\label{ch:xvii-topologia-odmian}

\section{Motywacja, obiekty i zależności logiczne}
Wcześniejsze rozdziały dostarczyły homologii singularnej, dualności
Poincarégo, klas charakterystycznych, struktur Hodge'a, teorii snopów
oraz monodromii. Teraz pytamy, jakie \emph{ograniczenia topologiczne}
narzuca pochodzenie przestrzeni z geometrii zespolonej algebraicznej.
W odróżnieniu od samej rozmaitości gładkiej, gładka odmiana projektowa
nosi wyróżnioną klasę kohomologii stopnia dwa, co prowadzi do
silnych zależności między jej grupami kohomologii i formami przecięcia.

Przez $X\subset\mathbb P^N_\C$ rozumiemy odtąd niepustą,
gładką, \emph{spójną} zespoloną odmianę projektową wymiaru zespolonego
$n$. Rozważamy jej analityfikację jako zwartą rozmaitość
rzeczywistą wymiaru $2n$, z orientacją zespoloną. Współczynniki
całkowite oznaczamy jawnie przez $\Z$; dla rozkładu Hodge'a używamy
$\C$, a w twierdzeniu o silnym Lefschetzu początkowo $\Q$ lub $\R$.
Nie utożsamiamy odmian osobliwych z rozmaitościami gładkimi: dla
pierwszych właściwym narzędziem bywa homologia przecięć.

\begin{definicja}[Klasa hiperpłaszczyzny i operator Lefschetza]
Niech $\mathcal O_X(1)$ oznacza ograniczenie wiązki hiperpłaszczyznowej
z $\mathbb P^N_\C$ do $X$, a $h=c_1(\mathcal O_X(1))\in
H^2(X;\Z)$. Niech $\omega$ będzie formą Kählera, której klasa
de Rhama jest obrazem $h$ w $H^2(X;\R)$ (po przyjęciu stosownej
normalizacji krzywizny). Definiujemy
$L:H^k(X;\Q)\to H^{k+2}(X;\Q)$ przez $L(\alpha)=h\smile\alpha$;
na kohomologii rzeczywistej jest to mnożenie przez $[\omega]$.
\end{definicja}
\begin{uwaga}
Gdy $h$ ma znaczenie jako całkowita klasa Cherna, pozostawiamy
zapis $h$. Gdy odwołujemy się do analizy Kählerowskiej,
używamy gładkiej formy $\omega$. Przejście między nimi wymaga
porównania kohomologii singularnej i de Rhama, przedstawionego
wcześniej w skrypcie.
\end{uwaga}

\section{Słabe twierdzenie Lefschetza o hiperpłaszczyźnie}
Niech $H\subset\mathbb P^N_\C$ będzie hiperpłaszczyzną taką,
że $Y=X\cap H$ jest gładkie i niepuste. Z twierdzenia Bertiniego
wiemy, że taki $H$ można wybrać z niepustego zbioru Zariski-otwartego
w przestrzeni dualnej (dla $n\ge1$ i odpowiedniego osadzenia).

% Moduł wspólny, wykorzystywany w geometrii i topologii algebraicznej.
% Dwa historyczne identyfikatory odsyłaczy odpowiadają jednemu twierdzeniu.
% Guard zapobiega podwójnemu składowi w książce łączącej oba rozdziały.
\ifcsname biblioteka@slabelefschetz@loaded\endcsname
Twierdzenie wraz z dowodem zostało już sformułowane wcześniej
(w twierdzeniu~\ref{tw:ga10-weak}).
\else
\expandafter\gdef\csname biblioteka@slabelefschetz@loaded\endcsname{}
\begin{twierdzenie}[Słabe twierdzenie Lefschetza]
\label{tw:ga10-weak}\label{thm:xvii-weak-lefschetz}
Niech $X\hookrightarrow\mathbb P^N_{\mathbb C}$ będzie gładką,
spójną odmianą projektową nad $\mathbb C$ o wymiarze zespolonym
$n\ge 2$. Niech $H\subset\mathbb P^N_{\mathbb C}$ będzie hiperpłaszczyzną,
a $Y=X\cap H$ niepustym gładkim przekrojem, z zanurzeniem
$i\colon Y\hookrightarrow X$. Wówczas homomorfizm kohomologii singularnej
\[
 i^*\colon H^q(X^{\mathrm{an}};\mathbb Z)
 \longrightarrow H^q(Y^{\mathrm{an}};\mathbb Z)
\]
jest izomorfizmem dla $q<n-1$ oraz monomorfizmem dla $q=n-1$.
Wersja kohomologiczna z dowolnymi współczynnikami w ciele
wynika z tej samej konstrukcji. W szczególności $Y$ jest spójna. Dla odpowiednio ustalonych
punktów bazowych indukowane odwzorowanie
$\pi_j(Y^{\mathrm{an}})\to\pi_j(X^{\mathrm{an}})$
jest izomorfizmem dla $1\le j<n-1$ oraz epimorfizmem dla $j=n-1$.
\end{twierdzenie}
\begin{proof}[Redukcja do geometrycznego twierdzenia Lefschetza]
Korzystamy z następującego \emph{głębokiego wyniku zewnętrznego}:
przestrzeń $X^{\mathrm{an}}$ ma typ homotopijny przestrzeni
otrzymanej z $Y^{\mathrm{an}}$ przez dołączenie komórek
rzeczywistego wymiaru co najmniej $n$. Dowód tego wyniku
wymaga zespolonej teorii Morse'a i funkcji plurisubharmonicznych
(lub odpowiedniej teorii znikania). Nie wyprowadzamy go z samej
gładkości $X$.

Względny kompleks komórkowy pary $(X^{\mathrm{an}},Y^{\mathrm{an}})$
nie ma komórek w stopniach mniejszych niż $n$, więc
$H_q(X^{\mathrm{an}},Y^{\mathrm{an}};\mathbb Z)=0$ dla $q<n$.
Względna aproksymacja komórkowa daje również
$\pi_j(X^{\mathrm{an}},Y^{\mathrm{an}})=0$ dla $j<n$.
Długi ciąg dokładny homotopii pary zapewnia stąd podane
izomorfizmy i epimorfizm. Należy przy tym pamiętać, że
w niskich stopniach grupy homotopii wymagają właściwej
interpretacji z ustalonymi punktami bazowymi.

Analogicznie $H^q(X^{\mathrm{an}},Y^{\mathrm{an}};\mathbb Z)=0$
dla $q<n$. W ciągu dokładnym kohomologii pary
\[
 \cdots\longrightarrow H^q(X,Y;\mathbb Z)
 \longrightarrow H^q(X;\mathbb Z)
 \xrightarrow{i^*} H^q(Y;\mathbb Z)
 \longrightarrow H^{q+1}(X,Y;\mathbb Z)\longrightarrow\cdots
\]
(przez $X,Y$ w tym ciągu rozumiemy przestrzenie analityczne)
dla $q<n-1$ znikają obie grupy względne sąsiadujące z $i^*$,
a dla $q=n-1$ znika grupa z lewej. Otrzymujemy odpowiednio
izomorfizm i monomorfizm. Ponieważ $n\ge2$, dla $q=0$
otrzymujemy izomorfizm $H^0(X;\mathbb Z)\simeq H^0(Y;\mathbb Z)$,
a spójność $X$ daje spójność $Y$.
\end{proof}
\fi


\begin{przyklad}[Przekrój powierzchni projektowej]
Dla $X=\mathbb P^2_\C$ i ogólnej prostej projektowej
$Y\cong\mathbb P^1_\C$ mamy $n=2$.
Słabe twierdzenie zapewnia izomorfizm w stopniu $0$ oraz
injektywność $H^1(X;\Z)\hookrightarrow H^1(Y;\Z)$.
\emph{Nie} twierdzi natomiast, że $H^2(X;\Z)\to H^2(Y;\Z)$
jest izomorfizmem, choć w tym konkretnym przykładzie oba
pierścienie mają naturalne generatory w stopniu dwa i restrykcja
jest izomorfizmem. To rozróżnienie zapobiega błędnemu
poszerzaniu zakresu twierdzenia.
\end{przyklad}

\begin{przyklad}[Krzywa jako przekrój powierzchni]
Jeśli $X\subset\mathbb P^3_\C$ jest gładką powierzchnią,
a $Y=X\cap H$ gładką krzywą, to $\pi_1(Y)\twoheadrightarrow\pi_1(X)$.
Nie oczekujemy na ogół izomorfizmu: krzywa może mieć
znacznie bogatszą grupę podstawową niż powierzchnia.
\end{przyklad}

\section{Silne twierdzenie Lefschetza i reprezentacje $\mathfrak{sl}_2$}
\begin{twierdzenie}[Silne twierdzenie Lefschetza]
\label{thm:xvii-hard-lefschetz}
Niech $X$ będzie gładką spójną zespoloną odmianą projektową
wymiaru $n$, a $h=c_1(\mathcal O_X(1))$.
Dla każdego $0\le k\le n$ mnożenie przez $h^{n-k}$ daje izomorfizm
\[
L^{n-k}:H^k(X;\Q)\xrightarrow{\ \simeq\ }H^{2n-k}(X;\Q).
\]
To samo zachodzi na kohomologii rzeczywistej dla każdej
klasy Kählerowskiej $[\omega]$.
\end{twierdzenie}

\begin{lemat}[Rachunek łańcuchów Lefschetza]
\label{lem:xvii-sl2}
Niech $V=\bigoplus_{k=0}^{2n}V^k$ będzie skończeniewymiarową
przestrzenią z dodatnio określonym iloczynem skalarnym,
$L:V^k\to V^{k+2}$, a $\Lambda=L^*$.
Załóżmy, że $[\Lambda,L]|_{V^k}=(n-k)\id$.
Wtedy dla $0\le k\le n$ operator
$L^{n-k}:V^k\to V^{2n-k}$ jest izomorfizmem.
\end{lemat}
\begin{proof}
Wprowadźmy $P^k=\ker(\Lambda:V^k\to V^{k-2})$.
Jeżeli $v\in P^k$ i $r\ge1$, to indukcja z komutatora daje
\[
\Lambda L^rv=r(n-k-r+1)L^{r-1}v.
\]
Ponieważ $\Lambda=L^*$, dla $v\ne0$ mamy
$\|L^rv\|^2=r(n-k-r+1)\|L^{r-1}v\|^2$.
Dla $r\le n-k$ współczynnik jest dodatni; przy
$r=n-k+1$ jest zerowy, zatem $L^{n-k+1}v=0$.
Ponadto każdy stopień dopuszcza prostopadły rozkład
$V^k=P^k\oplus LV^{k-2}$, gdyż $(\operatorname{im}L)^\perp=\ker\Lambda$.
Powtarzając rozkład (oraz korzystając z tożsamości komutatorów
w każdym stopniu), dostajemy sumę prostą
$V^k=\bigoplus_{r\ge0}L^rP^{k-2r}$, przy czym
$P^m=0$ dla $m>n$: gdy $v\in P^m$ i $m>n$, to
$\|Lv\|^2=(n-m)\|v\|^2<0$ dla $v\ne0$, sprzeczność.
Każdy łańcuch $v,Lv,\ldots,L^{n-m}v$ ma dokładnie
jeden niezerowy wektor w stopniach $m+2r$,
$0\le r\le n-m$. Obracając go działaniem $L^{n-k}$
na stopniu $k$ otrzymujemy bijekcję pomiędzy
odpowiednimi składnikami w stopniach $k$ i $2n-k$.
To dowodzi twierdzenia.
\end{proof}
\begin{proof}[Dowód silnego Lefschetza przy użyciu tożsamości Kählerowskich]
Rozważmy rzeczywistą przestrzeń harmonicznych form
$V^k=\mathcal H^k(X)$ dla ustalonej metryki Kählerowskiej.
Operator $L\alpha=\omega\wedge\alpha$ zachowuje formy
harmoniczne, podobnie jak jego formalny adjunkt $\Lambda$.
Istotnym analitycznym faktem są \emph{tożsamości Kählerowskie}
$[\Delta,L]=[\Delta,\Lambda]=0$ oraz
$[\Lambda,L]|_{\mathcal H^k}=(n-k)\id$.
Pełne wyprowadzenie ich w lokalnej ramce unitarnej należy do
analitycznej teorii Kählera; w tym rozdziale traktujemy je jako
jawnie przyjęty lemat geometryczny.
Lemat~\ref{lem:xvii-sl2} daje teraz izomorfizm
$L^{n-k}:\mathcal H^k(X)\to\mathcal H^{2n-k}(X)$.
Twierdzenie Hodge'a o reprezentantach harmonicznych
przenosi go na $H^k(X;\R)\to H^{2n-k}(X;\R)$.
Jeżeli $[\omega]=h$ jest klasą całkowitą,
macierz $L^{n-k}$ w bazach wymiernych ma współczynniki
w $\Q$. Jej wyznacznik nie znika nad $\R$, więc nie znika
nad $\Q$, co daje izomorfizm na kohomologii wymiernej.
\end{proof}

\begin{wniosek}[Symetria liczb Bettiego]
Dla gładkiej spójnej odmiany projektowej $X$ wymiaru $n$
$\dim_\Q H^k(X;\Q)=\dim_\Q H^{2n-k}(X;\Q)$.
Ponadto $b_{2j+1}(X)$ jest liczbą parzystą.
\end{wniosek}
\begin{proof}
Pierwsza równość wynika z silnego twierdzenia Lefschetza
(dla $k\le n$) i symetrii.
Z rozkładu Hodge'a $H^{2j+1}(X;\C)=
\bigoplus_{p+q=2j+1}H^{p,q}(X)$ i sprzężenia
$H^{p,q}=\overline{H^{q,p}}$ wynika, że suma wymiarów
jest parzysta, ponieważ $p\ne q$ przy $p+q$ nieparzystym.
\end{proof}

\section{Prymitywna kohomologia i relacje Hodge'a--Riemanna}
\begin{definicja}[Kohomologia prymitywna]
Dla $0\le k\le n$ definiujemy
$P^k(X;\C)=\ker(L^{n-k+1}:H^k(X;\C)\to H^{2n-k+2}(X;\C))$.
Korzystając z tożsamości Kählerowskich, można równoważnie
charakteryzować te klasy przez harmoniczne reprezentanty
$\alpha$ spełniające $\Lambda\alpha=0$.
Dla $p+q=k$ piszemy $P^{p,q}=P^k(X;\C)\cap H^{p,q}(X)$.
\end{definicja}
\begin{twierdzenie}[Rozkład Lefschetza]
\label{thm:xvii-decomposition}
Dla $0\le k\le n$ mamy rozkład
$H^k(X;\C)=\bigoplus_{r\ge0}L^rP^{k-2r}(X;\C)$,
przy czym składniki z ujemnym indeksem są zerowe.
\end{twierdzenie}
\begin{proof}
Lemat~\ref{lem:xvii-sl2} zastosowany do harmonicznych form
dostarcza rozkładu ortogonalnego na łańcuchy pierwotne.
Ponieważ $\omega$ jest typu $(1,1)$, mnożenie przez
$\omega$ zachowuje dwustopniowość Hodge'a, więc
rozłożenie jest zgodne z bigradowaniem.
\end{proof}
\begin{twierdzenie}[Relacje Hodge'a--Riemanna]
\label{thm:xvii-HR}
Dla $X$ i $\omega$ jak wyżej, $k=p+q\le n$ oraz
niezerowej klasy $\alpha\in P^{p,q}$ zachodzi
\[
(-1)^{k(k-1)/2}\,\mathrm i^{p-q}
\int_X\alpha\wedge\overline\alpha\wedge\omega^{n-k}>0.
\]
Całkę rozumiemy przez harmoniczne reprezentanty,
a wynik po lewej jest rzeczywistą liczbą dodatnią.
\end{twierdzenie}
\begin{proof}[Dodatniość z tożsamości gwiazdki Hodge'a]
Klasa prymitywna posiada harmonicznego reprezentanta
prymitywnego $a$ typu $(p,q)$. W ortonormalnej ramce
unitarnej dla metryki Kählerowskiej zachodzi punktowa
tożsamość
\[
*\overline a=
\frac{(-1)^{k(k-1)/2}\mathrm i^{p-q}}{(n-k)!}
\overline a\wedge\omega^{n-k},
\]
przy konwencji $\omega=\mathrm i\sum_j\theta^j\wedge
\overline{\theta^j}$ (z odpowiednio znormalizowaną ramką).
Tożsamość uzasadnia się sprawdzeniem na monomialnych
formach prymitywnych i liniowością; jest to potrzebny
lemat liniowej algebry hermitowskiej. Stąd
$(-1)^{k(k-1)/2}\mathrm i^{p-q}\int_X a\wedge
\overline a\wedge\omega^{n-k}=(n-k)!\int_X a\wedge *\overline a>0$.
Ostatnia całka jest normą $L^2$ niezerowej formy.
\end{proof}
\begin{uwaga}[Ścisły zakres dowodu]
Tożsamości Kählerowskie oraz punktową tożsamość dla
prymitywnej formy przyjmujemy jako lematy analitycznej
teorii Kählera; powyższy argument wyprowadza z nich
relacje i znaki. Nie jest to samodzielny dowód całej
teorii Hodge'a i twierdzenia Hodge'a--Riemanna od podstaw.
\end{uwaga}

\section{Formy przecięcia i twierdzenie o sygnaturze powierzchni}
\begin{definicja}[Forma przecięcia]
Dla gładkiej spójnej zespolonej powierzchni projektowej $S$
(przez \emph{powierzchnię} rozumiemy tutaj wymiar zespolony $2$)
definiujemy na $H^2(S;\R)$ formę
$Q_S(\alpha,\beta)=\int_S\alpha\wedge\beta$.
Z dualności Poincarégo jest ona niedegenerowana;
na $H^2(S;\Z)/\operatorname{Tor}$ wyznacza całkowitą
formę unimodularną.
\end{definicja}
\begin{twierdzenie}[Twierdzenie Hodge'a o indeksie]
\label{thm:xvii-hodge-index}
Jeśli $S$ jest gładką, spójną zespoloną powierzchnią projektową,
to sygnatura rzeczywistej formy $Q_S$ wynosi
\[
\bigl(b_2^+(S),b_2^-(S)\bigr)
=\bigl(1+2h^{2,0}(S),\ h^{1,1}(S)-1\bigr),
\]
a zatem $\sigma(S)=2+2h^{2,0}(S)-h^{1,1}(S)$.
W szczególności przecięcie na rzeczywistej przestrzeni
klas $(1,1)$ ma jeden kierunek dodatni, a na części
ortogonalnej do klasy Kählerowskiej jest ujemnie określone.
\end{twierdzenie}
\begin{proof}
Rozkład Hodge'a daje
$H^2(S;\C)=H^{2,0}\oplus H^{1,1}\oplus H^{0,2}$.
Przez $L$ oznaczamy mnożenie przez $[\omega]$.
Dekompozycja Lefschetza w stopniu $2$ przyjmuje postać
$H^2(S;\R)=\R[\omega]\oplus P^2(S;\R)$.
Ponieważ $\int_S\omega^2>0$, pierwsza prosta jest dodatnia.
Dla $a\in P^{1,1}\cap H^2(S;\R)$ relacja
Hodge'a--Riemanna w stopniu $k=2$ mówi
$-\int_S a\wedge a>0$ dla $a\ne0$.
Dla niezerowej klasy $b\in H^{2,0}$ mamy
$\int_S b\wedge\overline b>0$ (przypadek $k=2$, $p=2,q=0$).
Dodatniość na rzeczywistej płaszczyźnie
rozpiętej przez $\Re b,\Im b$ wynika stąd, że
$\int_S b\wedge b=0$ z powodów typu, a
$Q(\Re b,\Re b)=Q(\Im b,\Im b)
=\frac12\int_S b\wedge\overline b>0$ i składniki
mieszane są zerowe. Z powodów typu część $H^{1,1}$ jest ortogonalna do
$H^{2,0}\oplus H^{0,2}$ względem $Q$, natomiast
$H^{2,0}$ i $H^{0,2}$ \emph{nie} są wzajemnie ortogonalne:
parują ze sobą w całce $\int_S b\wedge\overline b$.
Na rzeczywistej przestrzeni $(H^{2,0}\oplus H^{0,2})\cap
H^2(S;\R)$ forma $Q$ jest dodatnio określona,
ponieważ każdy jej element można zapisać jako $b+\overline b$
i $Q(b+\overline b,b+\overline b)=2\int_S b\wedge\overline b>0$
dla $b\ne0$.
Liczba dodatnich kierunków to $1+2h^{2,0}$,
a ujemnych $h^{1,1}-1$, co kończy dowód.
\end{proof}
\begin{przyklad}[Płaszczyzna projektowa]
Dla $S=\mathbb P^2_\C$ mamy
$H^*(S;\Z)=\Z[h]/(h^3)$, $\deg h=2$,
$\int_S h^2=1$, stąd forma przecięcia na
$H^2(S;\Z)$ jest macierzą $(1)$ i $\sigma(S)=1$.
W tym przykładzie $h^{2,0}=0$ i $h^{1,1}=1$.
\end{przyklad}

\begin{przyklad}[Iloczyn dwóch prostych projektowych]
Niech $S=\mathbb P^1_\C\times\mathbb P^1_\C$.
Dla $a=\mathrm{pr}_1^*h$ i $b=\mathrm{pr}_2^*h$ mamy
$H^2(S;\Z)=\Z a\oplus\Z b$, $a^2=b^2=0$ oraz
$\int_S a\wedge b=1$ (formuła Künnetha).
Macierz przecięcia w bazie $(a,b)$ jest więc
$\begin{psmallmatrix}0&1\\1&0\end{psmallmatrix}$,
a jej wartości własne wynoszą $1$ i $-1$.
Ponieważ $h^{2,0}(S)=0$, $h^{1,1}(S)=2$,
twierdzenie o indeksie Hodge'a daje ten sam podpis $(1,1)$.
\end{przyklad}

\section{Rozdmuchanie w punkcie: geometria i topologia}
\begin{definicja}[Rozdmuchanie w punkcie]
Niech $S$ będzie gładką zespoloną powierzchnią i $p\in S$.
Rozdmuchaniem $\pi:\widetilde S=\operatorname{Bl}_pS\to S$
nazywamy odmianę, która jest izomorficzna z $S\setminus\{p\}$
poza dzielnikiem wyjątkowym $E=\pi^{-1}(p)\cong\mathbb P^1_\C$.
Lokalnie nad kartą ze współrzędnymi $(x,y)$ wokół $p$ jest to
pododmiana $\{((x,y),[s:t]):xt=ys\}$ w
$\C^2\times\mathbb P^1$ z naturalnym rzutowaniem.
\end{definicja}
% Wspólna podstawa geometryczna dla rozdmuchania punktu.
% Wariant opisu przez normalną i przez konormalną: obydwa zachowane.
% Historyczne etykiety występują tylko przy pierwszym włączeniu.
\ifcsname biblioteka@rozdmuchanie@normalna@loaded\endcsname
Wiązkę normalną dzielnika wyjątkowego obliczono wcześniej
w lemacie~\ref{lem:xvii-exceptional}.
\else
\expandafter\gdef\csname biblioteka@rozdmuchanie@normalna@loaded\endcsname{}
\begin{lemat}[Wiązka normalna dzielnika wyjątkowego]
\label{st:ga6-exceptional-normal}\label{lem:xvii-exceptional}
Niech $S$ będzie gładką zespoloną powierzchnią, $p\in S$, a
$\pi:\widetilde S=\operatorname{Bl}_pS\to S$ jej rozdmuchaniem
w punkcie $p$. Wtedy $E=\pi^{-1}(p)\simeq\mathbb P^1_{\mathbb C}$,
$N_{E/\widetilde S}\simeq\mathcal O_{\mathbb P^1}(-1)$ i
$E\cdot E=\deg N_{E/\widetilde S}=-1$.
Ostatnią liczbę rozumiemy jako lokalne samoprzecięcie krzywej,
nawet jeśli $S$ nie jest zwarta.
\end{lemat}
\begin{proof}
Wybierzmy współrzędne $(x,y)$ powierzchni $S$ z początkiem w $p$.
Rozdmuchanie jest lokalnie pododmianą
$\{((x,y),[s:t])\colon xt=ys\}\subset\mathbb C^2\times\mathbb P^1$.
W mapie $s\ne0$ piszemy $v=t/s$, $y=xv$, a w mapie $t\ne0$
piszemy $u=s/t$, $x=uy$. Na części wspólnej zachodzi
$u=v^{-1}$ i $y=vx$. W pierwszej mapie ideał $E$ jest
generowany przez $x$, w drugiej przez $y$.

\emph{Opis przez wiązkę normalną.}
W punkcie $[1:v]\in E$ kierunek normalny do $E$ jest
reprezentowany przez wektor $(x,xv)=x(1,v)$.
Po zmianie mapy wektory te wyznaczają dokładnie
podwiązki prostych rozpiętych przez $(s,t)$ w
$\mathbb P^1\times\mathbb C^2$; jest to tautologiczna
wiązka $\mathcal O_{\mathbb P^1}(-1)$.

\emph{Równoważny opis przez wiązkę konormalną.}
Lokalne generatory $x$ i $y$ wiązki
$\mathcal I_E/\mathcal I_E^2$ spełniają $y=vx$;
odpowiada to wiązce $\mathcal O_{\mathbb P^1}(1)$.
Jej dualna jest więc $N_{E/\widetilde S}\cong\mathcal O(-1)$.
Ponieważ stopień $\mathcal O_{\mathbb P^1}(-1)$ wynosi $-1$,
formuła samoprzecięcia
$E^2=\langle c_1(N_{E/\widetilde S}),[E]\rangle$
daje $E^2=-1$.
\end{proof}
\fi

% Wspólna część topologiczna: kohomologia, forma przecięć,
% liczby Bettiego, charakterystyka Eulera i sygnatura.
% Dwie historyczne etykiety identyfikują jedno twierdzenie.
\ifcsname biblioteka@rozdmuchanie@kohomologia@loaded\endcsname
Kohomologię rozdmuchania punktu obliczono wcześniej
w twierdzeniu~\ref{thm:xvii-blowup}.
\else
\expandafter\gdef\csname biblioteka@rozdmuchanie@kohomologia@loaded\endcsname{}
\begin{twierdzenie}[Kohomologia i forma przecięć rozdmuchania punktu]
\label{tw:ga10-blowup}\label{thm:xvii-blowup}
Niech $S$ będzie gładką, zwartą, spójną zespoloną powierzchnią,
$p\in S$, $\pi:\widetilde S=\operatorname{Bl}_pS\to S$ rozdmuchaniem,
$E=\pi^{-1}(p)$ oraz $e=\operatorname{PD}_{\widetilde S}[E]\in
H^2(\widetilde S;\mathbb Z)$ (symbol $\operatorname{PD}$ oznacza dualność
Poincarégo). Wówczas
\[
H^2(\widetilde S;\mathbb Z)=\pi^*H^2(S;\mathbb Z)\oplus\mathbb Z e.
\]
Dla każdej $\alpha\in H^2(S;\mathbb Z)$ zachodzi
$\langle \pi^*\alpha\smile e,[\widetilde S]\rangle=0$ i
$\langle e^2,[\widetilde S]\rangle=-1$.
W innych stopniach $\pi^*:H^k(S;\mathbb Z)\to
H^k(\widetilde S;\mathbb Z)$ jest izomorfizmem.
Przez $\sigma(M)$ rozumiemy sygnaturę rzeczywistej formy
przecięć na $H^2(M;\mathbb R)$; $b_2$ oznacza drugą liczbę
Bettiego, a $\chi_{\mathrm{top}}$ charakterystykę Eulera.
W szczególności
\[
 b_2(\widetilde S)=b_2(S)+1,\qquad
 \chi_{\mathrm{top}}(\widetilde S)=\chi_{\mathrm{top}}(S)+1,
 \qquad \sigma(\widetilde S)=\sigma(S)-1.
\]
\end{twierdzenie}
\begin{proof}
Wybierzmy małą zamkniętą kulę $B\subset S$ o środku $p$,
$U=S\setminus\operatorname{int}B$, i niech
$\widetilde B=\pi^{-1}(B)$. Przestrzeń $\widetilde B$ jest
otoczeniem dyskowym $E\simeq\mathbb P^1(\mathbb C)\simeq S^2$
i deformacyjnie retraktuje się na $E$. Brzegi $B$ oraz
$\widetilde B$ są identyfikowane przez $\pi$ ze sferą $S^3$.
Zatem $S=U\cup B$ i $\widetilde S=U\cup\widetilde B$
mają na przecięciach przestrzeń homotopijnie równoważną $S^3$.

Z ciągów Mayera--Vietorisa (ze współczynnikami $\mathbb Z$)
i równości $H^1(S^3)=H^2(S^3)=0$ dostajemy
\[
 H^2(S;\mathbb Z)\simeq H^2(U;\mathbb Z),\qquad
 H^2(\widetilde S;\mathbb Z)
 \simeq H^2(U;\mathbb Z)\oplus H^2(E;\mathbb Z).
\]
Pierwsza składowa ostatniego izomorfizmu jest obrazem $\pi^*$,
ponieważ $\pi|_{\widetilde B}$ ma obraz w ściągalnej kuli $B$.
Ograniczenie $e$ do $E$ jest klasą
$c_1(N_{E/\widetilde S})=-\eta$, gdzie
$\eta$ jest dodatnim generatorem $H^2(E;\mathbb Z)$;
wynika to z lematu~\ref{lem:xvii-exceptional}.
Dlatego $e$ generuje dodatkowy składnik całkowity.

Dla jasności sprawdzamy także stopień trzeci.
Z dualności Poincarégo--Lefschetza lub z długiego ciągu pary
$(U,\partial U)$ wynika, że homomorfizm łączący
$H^3(\partial U;\mathbb Z)\to H^4(U,\partial U;\mathbb Z)$
jest izomorfizmem: generator sfery granicznej przechodzi
na względną klasę orientacji spójnej rozmaitości $U$.
Stąd $H^3(U;\mathbb Z)\to H^3(\partial U;\mathbb Z)$
jest zerowy. Porównanie obu ciągów Mayera--Vietorisa
pokazuje, że $\pi^*$ jest izomorfizmem również w stopniu trzecim.
Stopnie zero i jeden porównujemy bezpośrednio z ciągów,
a w stopniu czwartym $\pi$ ma stopień $1$, więc
$\pi^*:H^4(S;\mathbb Z)\to H^4(\widetilde S;\mathbb Z)$
jest izomorfizmem (obie grupy są izomorficzne z $\mathbb Z$).

Ortogonalność wynika ze wzoru projekcji i tego, że
$\pi|_E$ jest odwzorowaniem stałym:
\[
 \langle\pi^*\alpha\smile e,[\widetilde S]\rangle
 =\langle (\pi|_E)^*\alpha,[E]\rangle=0.
\]
Samoprzecięcie $e^2=-1$ daje wspomniany lemat.
Ponieważ $\pi$ ma stopień $1$, forma przecięć na
$\pi^*H^2(S;\mathbb R)$ jest identyczna z formą na $S$.
Otrzymujemy więc rozkład ortogonalny
$Q_{\widetilde S}\cong Q_S\oplus(-1)$ na wolnych częściach
kohomologii. To daje $\sigma(\widetilde S)=\sigma(S)-1$.
Zmiana wyłącznie drugiej liczby Bettiego powoduje
$\chi_{\mathrm{top}}(\widetilde S)=\chi_{\mathrm{top}}(S)+1$.
\end{proof}
\fi

% Wniosek o sygnaturze i charakterystyce Eulera zawiera teraz twierdzenie wspólne.
\begin{przyklad}[Rozdmuchanie $\mathbb P^2$]
Dla $S=\mathbb P^2_\C$ i
$\widetilde S=\operatorname{Bl}_p\mathbb P^2$ otrzymujemy bazę
$(H,E)$ dla $H^2(\widetilde S;\Z)$ modulo torsję,
gdzie $H=\pi^*h$, $E=e$, i macierz przecięcia
\[
\begin{pmatrix}1&0\\0&-1\end{pmatrix}.
\]
Zatem $b_2=2$, $\sigma=0$, $\chi=4$.
Klasa $H$ jest całkowita, ale nie jest klasą Kählerowską
na rozdmuchaniu: jej całka po krzywej wyjątkowej wynosi zero.
Dla odpowiednio małego $\varepsilon>0$ klasa
$H-\varepsilon E$ jest Kählerowska.
\end{przyklad}

\section{Schemat geometryczny rozdmuchania}
\begin{center}
\begin{tikzpicture}[>=Stealth,scale=1,
  every node/.style={font=\small}]
  \draw[thick,rounded corners=8pt] (0,0) rectangle (4,2.3);
  \fill (2,1.13) circle (2pt);
  \node[below] at (2,1.0) {$p$};
  \node at (2,2.58) {$S$ (lokalnie)};
  \draw[thick,rounded corners=8pt] (7,0) rectangle (11,2.3);
  \draw[thick,manifoldblue] (9,1.28) ellipse (0.86 and 0.38);
  \draw[thick,manifoldblue,dashed] (8.14,1.28) arc (180:360:0.86 and 0.38);
  \node[below] at (9,0.63) {$E\cong\mathbb P^1$};
  \node at (9,2.58) {$\widetilde S$};
  \draw[->,very thick] (7,1.2) -- node[above] {$\pi$} (4,1.2);
  \node[font=\footnotesize,align=center] at (5.5,-0.55)
    {schemat: cały dzielnik $E$\ jest ściągany do $p$};
\end{tikzpicture}
\par\smallskip\small Rysunek poglądowy: $E$ jest zespoloną krzywą projektową,
a nie okręgiem w przestrzeni rzeczywistej.
\end{center}

\section{Dzielnik kanoniczny i charakterystyczne liczby Chernowskie}
\begin{twierdzenie}[Kanoniczna klasa rozdmuchania powierzchni]
\label{thm:xvii-canonical}
Jeżeli $S$ jest gładką zespoloną powierzchnią,
a $\pi:\widetilde S\to S$ jest rozdmuchaniem w punkcie,
to dla wiązek kanonicznych zachodzi
$K_{\widetilde S}\cong\pi^*K_S\otimes\mathcal O_{\widetilde S}(E)$.
Jeśli $S$ jest ponadto projektowa, to
$K_{\widetilde S}^2=K_S^2-1$.
\end{twierdzenie}
\begin{proof}
W karcie $y=ux$ różniczkowanie daje
$\dd x\wedge\dd y=x\,\dd x\wedge\dd u$.
Jacobian mapy $\pi$ ma zatem proste zero na $E$.
W drugiej karcie otrzymujemy zgodny rachunek, toteż
dzielnik ramifikacji ma krotność jeden wzdłuż $E$.
Formuła dla wiązki kanonicznej wynika z identyfikacji
lokalnych generatorów najwyższej potęgi wiązki kostycznej.
Po wzięciu pierwszej klasy Cherna mamy
$c_1(K_{\widetilde S})=\pi^*c_1(K_S)+e$.
Wymnażając kwadrat i korzystając z zerowania
$\pi^*c_1(K_S)\smile e$ oraz $e^2=-1$,
otrzymujemy $K_{\widetilde S}^2=K_S^2-1$.
\end{proof}

\section{Rozdmuchanie w gładkim centrum i formuła kohomologiczna}
\begin{twierdzenie}[Wzór na kohomologię rozdmuchania]
\label{thm:xvii-blowup-general}
Niech $X$ będzie gładką zespoloną odmianą projektową,
a $Z\subset X$ gładką domkniętą pododmianą o zespolonej
kodymensji $r\ge2$. Niech $\widetilde X=\operatorname{Bl}_Z X$,
$\pi:\widetilde X\to X$, $j:E\hookrightarrow\widetilde X$,
$p:E=\mathbb P(N_{Z/X})\to Z$ i
$\xi=c_1(\mathcal O_E(1))$.
Istnieje addytywny izomorfizm grup całkowitej kohomologii
\[
H^k(\widetilde X;\Z)\cong H^k(X;\Z)\oplus
\bigoplus_{a=1}^{r-1}H^{k-2a}(Z;\Z),
\]
realizowany przez $\pi^*$ oraz $j_*(\xi^{a-1}p^*(-))$,
gdzie $j_*$ jest homomorfizmem Gysina podnoszącym stopień o $2$.
Po przejściu na $\C$ rozkład jest zgodny ze strukturą Hodge'a,
przy czym nowy składnik $H^{p-a,q-a}(Z)$ przechodzi do
$H^{p,q}(\widetilde X)$.
\end{twierdzenie}
\begin{proof}[Wyprowadzenie z teorii par i formuły wiązki projektowej]
Mapy $\pi$ i $p$ dają diagram par $(\widetilde X,E)\to(X,Z)$.
Poza $E$ i $Z$ odwzorowanie jest izomorfizmem;
ponieważ $X$ i $\widetilde X$ są zwarte i Hausdorffa, odwzorowanie
$\pi$ indukuje homeomorfizm ilorazów
$\widetilde X/E\cong X/Z$. Istotnie, po ściągnięciu $E$ i $Z$
do jednego punktu dostajemy ciągłą bijekcję między przestrzeniami
zwartą i Hausdorffa; jest ona homeomorfizmem. Daje to
$H^k(X,Z;\Z)\cong H^k(\widetilde X,E;\Z)$.

Rozpatrzmy morfizm długich ciągów kohomologii par
$(X,Z)$ i $(\widetilde X,E)$; poziome mapy między grupami
względnymi są izomorfizmami, a na zamkniętych podprzestrzeniach
mamy $p^*:H^k(Z)\to H^k(E)$.
Twierdzenie Leraya--Hirscha dla $p:E\to Z$ zapewnia, że
$p^*$ jest monomorfizmem we wszystkich stopniach.
Lemat o wężu zastosowany do odpowiadających fragmentów
ciągów długich daje wtedy dokładny ciąg
\[
0\longrightarrow H^k(X;\Z)\xrightarrow{\pi^*}
H^k(\widetilde X;\Z)\xrightarrow{[j^*]}
H^k(E;\Z)/p^*H^k(Z;\Z)\longrightarrow0.
\]
W szczególności $\pi^*$ jest monomorfizmem, a $[j^*]$
oznacza restrykcję do $E$ po przejściu do ilorazu przez obraz
$p^*$. Dokładność tego ciągu, a nie jedynie formalna wzmianka
o pięciu homomorfizmach, jest potrzebnym krokiem porównania.

Z twierdzenia Leraya--Hirscha dla wiązki projektowej
$E=\mathbb P(N_{Z/X})\to Z$ mamy
$H^*(E;\Z)=\bigoplus_{a=0}^{r-1}\xi^a p^*H^{*-2a}(Z;\Z)$
addytywnie. Obraz $p^*$ jest jego składnikiem $a=0$.
Dla klas nowych na $\widetilde X$ wykorzystujemy
homomorfizm Gysina $j_*$ oraz tożsamość
$j^*j_*\beta=c_1(N_{E/\widetilde X})\smile\beta
=-\xi\smile\beta$.
W ilorazie $H^k(E)/p^*H^k(Z)$ klasy
$\xi^ap^*\beta$, $1\le a\le r-1$, tworzą dokładnie
pozostałe składniki rozkładu Leraya--Hirscha.
Stąd klasy $j_*(\xi^{a-1}p^*\beta)$ przechodzą przez
$[j^*]$ na $-\xi^ap^*\beta$. Zatem suma homomorfizmów
Gysina daje przekrój powyższego krótkiego ciągu dokładnego
(po zmianie znaku); wraz z $\pi^*$ otrzymujemy żądany
izomorfizm addytywny. Jest on określony przez jawne mapy.
Użytymi zewnętrznymi rezultatami pozostają formuła
Leraya--Hirscha dla wiązki projektowej, konstrukcja mapy
Gysina i wzór samoprzecięcia; nie ukrywamy tych zależności.
Zgodność typów Hodge'a wynika stąd, że $\pi^*,p^*$
są morfizmami struktur Hodge'a, $\xi$ ma typ $(1,1)$,
a $j_*$ podnosi oba indeksy o $1$.
\end{proof}
\begin{przyklad}[Rozdmuchanie w krzywej w trójwymiarowej odmianie]
Niech $X$ będzie gładką projektową trójwymiarową odmianą,
a $Z\subset X$ gładką spójną krzywą (zatem $r=2$).
Wtedy $H^3(\operatorname{Bl}_Z X;\Q)
\cong H^3(X;\Q)\oplus H^1(Z;\Q)$.
Jeśli $Z$ ma rodzaj $g$, to trzeci wymiar Bettiego rośnie o $2g$.
W szczególności operacja biracjonalna może zmienić środkową
kohomologię nawet wtedy, gdy nie zmienia dopełnienia centrum.
\end{przyklad}

\begin{przyklad}[Rozdmuchanie $\mathbb P^3$ wzdłuż gładkiej krzywej eliptycznej]
\label{prz:xvii-blowup-elliptic-audit}
Niech $X=\mathbb P^3_\C$, a $Z\subset\mathbb P^2_\C
\subset\mathbb P^3_\C$ będzie krzywą zadaną równaniem
$x^3+y^3+z^3=0$ w ustalonej płaszczyźnie projektowej.
Jej pochodne cząstkowe to $3x^2,3y^2,3z^2$;
nie mogą wszystkie znikać w punkcie projektowym.
Zatem $Z$ jest gładką krzywą płaską stopnia trzy.
Z formuły adjunkcji dla gładkiej krzywej płaskiej
otrzymujemy
\[
 K_Z\cong
 (K_{\mathbb P^2}\otimes\mathcal O_{\mathbb P^2}(3))|_Z
 \cong\mathcal O_Z.
\]
Ponieważ $\deg K_Z=2g(Z)-2$, mamy $g(Z)=1$.
Krzywa $Z$ jest więc eliptyczna (po wyborze punktu bazowego).

Jako pododmiana $X$ krzywa ma kodymensję zespoloną
$r=3-1=2$. Dla
$\widetilde X=\operatorname{Bl}_Z\mathbb P^3$
Twierdzenie~\ref{thm:xvii-blowup-general} redukuje się do
\[
 H^k(\widetilde X;\Z)\cong
 H^k(\mathbb P^3;\Z)\oplus H^{k-2}(Z;\Z).
\]
Pierwszy składnik jest równy $\Z$ w stopniach
$k=0,2,4,6$ i zero w pozostałych.
Dla drugiego używamy
$H^0(Z;\Z)\cong H^2(Z;\Z)\cong\Z$ oraz
$H^1(Z;\Z)\cong\Z^2$. Dostajemy wszystkie liczby Bettiego:
\[
 (b_0,b_1,b_2,b_3,b_4,b_5,b_6)
   = (1,0,2,2,2,0,1).
\]
Zatem
$\chi_{\mathrm{top}}(\widetilde X)
=1+2-2+2+1=4$.
Charakterystyka Eulera nie zmienia się względem
$\chi_{\mathrm{top}}(\mathbb P^3)=4$,
mimo że przybyły \emph{dwie} niezależne
klasy kohomologii stopnia trzeciego.

Wyjaśnijmy również ich typy Hodge'a.
Mamy $H^1(Z;\C)=H^{1,0}(Z)\oplus H^{0,1}(Z)$,
a każda z tych przestrzeni ma wymiar $1$.
Nowy składnik stopnia trzeciego jest obrazem
homomorfizmu Gysina
$j_*p^*:H^1(Z;\C)\to H^3(\widetilde X;\C)$.
Jako morfizm typu $(1,1)$ podnosi on
$(1,0)$ do $(2,1)$, a $(0,1)$ do $(1,2)$.
Stąd
\[
 h^{2,1}(\widetilde X)=h^{1,2}(\widetilde X)=1,
 \qquad h^{3,0}(\widetilde X)=h^{0,3}(\widetilde X)=0.
\]
Nowe klasy w stopniach drugim i czwartym
mają odpowiednio typy $(1,1)$ i $(2,2)$.
Ten rachunek pokazuje, że rozdmuchanie w centrum
dodatniego rodzaju może tworzyć środkową
kohomologię o nietrywialnych typach Hodge'a,
nawet gdy wyjściowa przestrzeń projektowa ma
jedynie klasy diagonalne.
\end{przyklad}

\begin{figure}[htbp]
\centering
\begin{tikzpicture}[>=Stealth,node distance=13mm,
  box/.style={draw,rounded corners=3pt,minimum width=27mm,
              minimum height=8mm,align=center,font=\small}]
\node[box] (a) {$H^2(Z)\cong\Z$};
\node[box,below=of a] (b) {$H^1(Z)\cong\Z^2$};
\node[box,below=of b] (c) {$H^0(Z)\cong\Z$};
\node[box,right=36mm of a] (aa) {nowe $H^4$};
\node[box,right=36mm of b] (bb) {nowe $H^3$};
\node[box,right=36mm of c] (cc) {nowe $H^2$};
\draw[->,thick,manifoldblue] (a) --
 node[above,font=\small] {$j_*p^*$} (aa);
\draw[->,thick,manifoldgreen!75!black] (b) --
 node[above,font=\small] {$j_*p^*$} (bb);
\draw[->,thick,manifoldred] (c) --
 node[above,font=\small] {$j_*p^*$} (cc);
\end{tikzpicture}
\caption{Przy rozdmuchaniu w centrum o kodymensji $2$
homomorfizm Gysina przesuwa stopnie kohomologii krzywej
o dwa. W szczególności $H^1(Z)$ tworzy nowe $H^3$ rozdmuchania.}
\label{fig:xvii-elliptic-blowup-audit}
\end{figure}

\section{Powiązanie z teorią snopów perwersyjnych i monodromią}
Twierdzenia z rozdziałów~\ref{ch:ih-perverse-xiii},
\ref{ch:xiv-hodge-monodromia} i~\ref{ch:xvi-szesc-operacji}
dają alternatywne podejście do hiperpłaszczyzn i rozdmuchań.
Dla właściwego morfizmu algebraicznego $f$ kompleksy $Rf_*IC_X$
można badać przez twierdzenie o rozkładzie.
Gdy $f$ jest rozdmuchaniem gładkiego centrum,
dodatkowe składniki kohomologiczne są geometrycznie
pochodnymi obrazami dzielnika wyjątkowego; ich przesunięcia
odpowiadają stopniom nowo powstających klas.
Nie należy jednak utożsamiać wzoru kohomologicznego
na rozdmuchanie z pełnym twierdzeniem BBDG:
pierwszy daje tu konkretną dekompozycję,
a drugie jest znacznie ogólniejszym twierdzeniem
w kategorii pochodnej snopów konstruktybilnych.

\section{Dalsze przykłady i kontrola intuicji}
\begin{przyklad}[Czemu dodatniość zmienia się po rozdmuchaniu?]
Na $\operatorname{Bl}_p\mathbb P^2$ forma przecięcia w bazie
$(H,E)$ ma sygnaturę $(1,1)$. Dla klasy rzeczywistej
$\kappa=H-\varepsilon E$ z $0<\varepsilon<1$ mamy
$\kappa^2=1-\varepsilon^2>0$ oraz
$\kappa\cdot E=\varepsilon>0$.
To jest spójne z możliwością reprezentowania $\kappa$
przez formę Kählera; jednak sama dodatniość kwadratu
i dodatnie przecięcie z $E$ nie wystarczają w ogólności
do wykazania Kählerowskości. W tym przykładzie
Kählerowskość dla małego $\varepsilon$ wynika z osobnego
twierdzenia o klasach Kählerowskich rozdmuchań.
\end{przyklad}
\begin{przyklad}[Rozdzielenie trzech twierdzeń Lefschetza]
\emph{Słabe twierdzenie Lefschetza} porównuje kohomologię $X$ i gładkiego
przekroju $Y=X\cap H$. \emph{Silne twierdzenie Lefschetza} mówi o
izomorfizmach $L^{n-k}$ wewnątrz kohomologii $X$.
\emph{Względne silne twierdzenie Lefschetza} dla właściwych map algebraicznych
jest twierdzeniem o perwersyjnych snopach kohomologii
$Rf_*IC_X$; nie należy go używać bez warunków
właściwości, projektowości i normalizacji perwersyjnej.
\end{przyklad}

\section{Bilans dowodów i przygotowanie do wspólnej książki}
W tym rozdziale samodzielnie przeprowadziliśmy topologiczne
konciągi dołączania komórek, liniowoalgebraiczny lemat
Lefschetza, dowód indeksu Hodge'a z relacji
Hodge'a--Riemanna oraz obliczenia rozdmuchania w punkcie.
\emph{Jawnymi twierdzeniami pomocniczymi} pozostają:
Bertini, geometryczne twierdzenie Lefschetza o dołączaniu
komórek, tożsamości Kählerowskie, twierdzenie Hodge'a
o reprezentantach harmonicznych, punktowa tożsamość
prymitywnych form Hodge'a oraz formuła wiązki projektowej.
Ich pełne dowody wymagają teorii, której nie zastępujemy
kilkoma zdaniami.

W późniejszym połączeniu z geometrią algebraiczną
warto zachować w tym skrypcie \emph{skutki topologiczne},
a konstrukcję snopów koherentnych, formułę wiązki
projektowej w przecięciach algebraicznych i pełne
teorie dzielników wyjaśnić w tomie algebraicznym.

